\begin{thema}{A}
\begin{erwthma}
    \item Να αποδείξετε ότι η εκθετική συνάρτηση $ f(x)=a^x $ με $ 0<a\neq 1 $ είναι παραγωγίσιμη στο $ \mathbb{R} $ με $ f'(x)=a^x\cdot\ln{a} $.
    \item Πότε λέμε ότι μια συνάρτηση $f$ στρέφει τα κοίλα προς τα άνω (ή είναι κυρτή) στο $\Delta$;
    \item Πότε μια συνάρτηση $f:A\to\mathbb{R}$ ονομάζεται συνάρτηση $1-1$;
    \item Να χαρακτηρίσετε τις προτάσεις που ακολουθούν, γράφοντας στο τετράδιό σας, δίπλα στο γράμμα που αντιστοιχεί σε κάθε πρόταση, τη λέξη \textbf{Σωστό}, αν η πρόταση είναι σωστή, ή \textbf{Λάθος}, αν η πρόταση είναι λανθασμένη.
    \begin{alist}
        \item Το πεδίο ορισμού της σύνθεσης $f \circ g$ ταυτίζεται πάντοτε με την τομή των πεδίων ορισμού της $f$ και της $g$.
        \item Η κατακόρυφη ευθεία $x=x_0$ λέγεται κατακόρυφη ασύμπτωτη αν έστω ένα πλευρικό όριο στο $x_0$ είναι $+\infty$ ή $-\infty$.
        \item Η $\hm x$ είναι κυρτή στο $(0, \pi)$.
        \item Αν ισχύει $F(x) = G(x) + c$, τότε η $F$ και η $G$ είναι παράγουσες της ίδιας συνάρτησης.
        \item Αν $\dint_a^\beta f(x)dx = 0$, τότε αναγκαστικά $a = \beta$.
    \end{alist}
    \item Παρατηρήστε την παρακάτω γραφική παράσταση της \textbf{παραγώγου} $f'(x)$ μιας συνεχούς συνάρτησης $f$.
    \begin{center}
    \begin{tikzpicture}
      \begin{axis}[
          aks_on,
          xmin=-3, xmax=4,
          ymin=-2, ymax=4,
          xtick={-1, 2}, xticklabels={$-1$, $2$},
          ytick=\empty,
          width=8cm, height=5cm
      ]
      \addplot[grafikh parastash, domain=-1.5:3, samples=100] {(x+1)*(x-2)};
      \node[anchor=south west] at (axis cs:2.5, 1) {$y = f'(x)$};
      \end{axis}
    \end{tikzpicture}
    \end{center}
    Σε ποιο από τα παρακάτω διαστήματα η αρχική συνάρτηση $f(x)$ είναι γνησίως φθίνουσα;
    \begin{tasks}[label=\let\textdexiakeraia\relax\alph*.](4)
        \task $[ -1, 2 ]$
        \task $(-\infty, -1]$
        \task $[2, +\infty)$
        \task $\mathbb{R}$
    \end{tasks}
\end{erwthma}
\end{thema}
